% ════════════════════════════════════════════
% 摘要
% ════════════════════════════════════════════
\begin{abstract}

社交平台上的帖子曝光量决定了内容能否被看见。本文以某化妆品品牌主页的 500 条帖子为研究对象，数据为帖级截面记录，含粉丝规模、帖子类型、内容主题、发布月份、星期、时段、付费标记 7 个输入与曝光总次数 1 个输出。分析要解决两个问题：一是判断哪些发帖决策真正改变了曝光，二是把回归系数翻译成可直接执行的投放建议。样本中粉丝规模与发布月份高度同步（Spearman 0.996、辅助回归 $R^2=0.9917$），两个变量同时进入模型会使系数不可识别，这构成了建模的首要障碍。

数据处理上，曝光量取 $\log(1+Y_1)$ 后偏度由 8.33 降至 0.73、峰度由 93.05 降至 0.82；唯一的缺失记录按众数填补为自然帖；名义与周期变量做独热编码，24 小时的发布时段按样本分布合并为凌晨、上午、下午晚间三组。共线判别采用三项判据：相关系数绝对值不低于 0.9、辅助回归解释度不低于 0.9、月内方差占比低于 5\%，三项同时满足即判定不可同时识别。据此剔除粉丝规模、保留月份固定效应，使最大方差膨胀因子由 143.0 降至 3.23。

以 $\log(1+Y_1)$ 为因变量估计五个候选模型，最终模型 $R^2=0.207$、十折交叉验证 $R^2=0.101$、交叉验证 RMSE 为 1.101，优于均值基线的 $-0.021$。残差的 Breusch--Pagan 检验 $p=0.0005$、Jarque--Bera 检验 $p<0.001$，推断统一采用 HC3 稳健标准误。Shapley 分解显示曝光方差的 37.1\% 来自帖子类型、28.0\% 来自月份，模型整体可解释 20.7\%，其余 79.3\% 来自未被观测的个体差异。

效应测算给出三条可执行结论：控制其他变量后付费帖子的曝光比自然帖子高 40.1\%（95\% 置信区间 10.8\%--77.2\%，代表性画像上相当于多 16,142 次）；视频帖是图片帖的 4.90 倍，但样本仅 7 条、区间为 2.23--10.79；主题 2 与主题 3 分别比主题 1 低 38.0\% 与 27.0\%。六种模型设定下付费提升幅度落在 37.4\%--49.6\%、视频倍数落在 3.52--5.41，方向与量级一致；按内容主题分组检验发现付费效应存在异质性（组内 $p=0.036$），效应主要集中在主题 1。上述结论适用于单品牌、帖级截面、对数尺度的场景，跨品牌推广或需要绝对曝光量预测时须重新标定。

\keywords{帖子曝光；对数变换；月份固定效应；HC3 稳健标准误；效应测算}

\end{abstract}
